\documentclass[10pt,a4paper]{amsart}
\usepackage[margin=19mm]{geometry}
\usepackage{amsmath,amssymb,mathtools}
\usepackage[scr=rsfs,cal=euler]{mathalfa}
\usepackage{xfrac}
\usepackage[unicode,hidelinks,bookmarks=false]{hyperref}
\newcommand{\ie}{i.e.\ }
\newcommand{\cf}{cf.\ }
\newcommand{\resp}{respectively\ }
\newcommand{\Subsection}{\subsection}
\providecommand{\nobreakdash}{\nobreak\hbox{-}}
\setlength{\emergencystretch}{2em}
\multlinegap0pt
% wrapper: the article is an imsart[aos] article.  That publisher class is not
% part of the pinned runtime; the AMS amsart wrapper below carries the pinned
% runtime preamble, the named packages the article itself loads, and then the
% the article's own macro and theorem declaration block copied verbatim from its own
% 3_preamble.sty (lines 61-184), so that every definition and every theorem-like
% environment keeps the article's own meaning, style and counter sharing.
\usepackage{amsthm}
\usepackage{xcolor}
\usepackage[capitalize,sort]{cleveref}
\newcommand{\blue}[1]{{\color{blue}#1}}
\newcommand{\red}[1]{{\color{red}#1}}
\newcommand{\gray}[1]{{\color{gray}#1}}
\newcommand{\violet}[1]{{\color{violet}#1}}

\newcommand{\commhb}[1]{{\textbf{{\violet{{[HB: #1]}}}}}}

%%%%%%%%%%%%%%%%%%%%%%%%%%%%%%%%%%%%%%%%%%%%%%
%%                                          %%
%% Uncomment next line to change            %%
%% the type of equation numbering           %%
%%                                          %%
%%%%%%%%%%%%%%%%%%%%%%%%%%%%%%%%%%%%%%%%%%%%%%
%\numberwithin{equation}{section}
%%%%%%%%%%%%%%%%%%%%%%%%%%%%%%%%%%%%%%%%%%%%%%
%%                                          %%
%% For Axiom, Claim, Corollary, Hypothezis, %%
%% Lemma, Theorem, Proposition              %%
%% use \theoremstyle{plain}                 %%
%%                                          %%
%%%%%%%%%%%%%%%%%%%%%%%%%%%%%%%%%%%%%%%%%%%%%%
\theoremstyle{plain}
\newtheorem{axiom}{Axiom}
\newtheorem{claim}[axiom]{Claim}
\newtheorem{theorem}{Theorem}[section]
\newtheorem{lemma}{Lemma}[section]
\newtheorem{corollary}[theorem]{Corollary}
\newtheorem{proposition}[theorem]{Proposition}

%%%%%%%%%%%%%%%%%%%%%%%%%%%%%%%%%%%%%%%%%%%%%%
%%                                          %%
%% For Assumption, Definition, Example,     %%
%% Notation, Property, Remark, Fact         %%
%% use \theoremstyle{remark}                %%
%%                                          %%
%%%%%%%%%%%%%%%%%%%%%%%%%%%%%%%%%%%%%%%%%%%%%%
\theoremstyle{remark}
\newtheorem{definition}[theorem]{Definition}
\newtheorem{remark}[theorem]{Remark}
\newtheorem*{example}{Example}
\newtheorem*{fact}{Fact}

%%%%%%%%%%%%%%%%%%%%%%%%%%%
%%% Algorithmic Command %%%
%%%%%%%%%%%%%%%%%%%%%%%%%%%

% \newcommand{\algorithmicinput}{\textbf{Input:}}
% \newcommand{\INPUT}{\item[\algorithmicinput]}
% \newcommand{\algorithmicoutput}{\textbf{Output:}}
% \newcommand{\OUTPUT}{\item[\algorithmicoutput]}
% \newcommand{\algorithmicinit}{\textbf{Initialization:}}
% \newcommand{\INIT}{\item[\algorithmicinit]}
% \newcommand{\algorithmiciter}{\textbf{Iteration:}}
% \newcommand{\ITER}{\item[\algorithmiciter]}

%%%%%%%%%%%%%%%%%%%%%%
%%% Important Sets %%%
%%%%%%%%%%%%%%%%%%%%%%
\newcommand{\reals}{\ensuremath{\mathbb{R}}}
\newcommand{\ints}{\ensuremath{\mathbb{Z}}}
\newcommand{\rats}{\ensuremath{\mathbb{Q}}}
\newcommand{\nats}{\ensuremath{\mathbb{N}}}
\newcommand{\comps}{\ensuremath{\mathbb{C}}}

%%%%%%%%%%%%%%%%%%%%%%%%
%%% Matrix Operators %%%
%%%%%%%%%%%%%%%%%%%%%%%%
\newcommand{\diag}{\operatorname{diag}}
\newcommand{\tr}{\operatorname{tr}}

%%%%%%%%%%%%%%%%%%%%%%%
%%% General Purpose %%%
%%%%%%%%%%%%%%%%%%%%%%%
\newcommand{\textand}{\enskip \text{and} \enskip}
\def\hat#1{\widehat{#1}}
\def\tilde#1{\widetilde{#1}}
\def\bar#1{\overline{#1}}

\newcommand{\sgn}{\text{sgn}}
\newcommand{\tsum}{{\textstyle\sum}}
\newcommand{\tprod}{{\textstyle\prod}}

\let\oldabs\abs % Store original \abs as \oldabs
\let\abs\undefined % "Undefine" \abs
\DeclarePairedDelimiter\abs{\lvert}{\rvert}

\let\oldnorm\norm   % <-- Store original \norm as \oldnorm
\let\norm\undefined % <-- "Undefine" \norm
\DeclarePairedDelimiter\norm{\lVert}{\rVert}
\DeclarePairedDelimiter\mnorm{\vvvert}{\vvvert}

\DeclarePairedDelimiter\inner{\langle}{\rangle}
\DeclarePairedDelimiter\ceil{\lceil}{\rceil}
\DeclarePairedDelimiter\floor{\lfloor}{\rfloor}

\DeclareMathOperator*{\argmin}{\arg\!\min}
\DeclareMathOperator*{\argmax}{\arg\!\max}

\newcommand{\Rho}{\mathrm{P}}
\newcommand{\vareps}{\varepsilon}
\newcommand{\cmpl}{\mathsf{C}}

%%%%%%%%%%%%%%%%%%%%%%%%%%%%%%%%%%
%%% Probability and Statistics %%%
%%%%%%%%%%%%%%%%%%%%%%%%%%%%%%%%%%
\renewcommand{\Pr}{\mathbb{P}}  % probability
\newcommand{\Ind}{\mathbb{I}}
\newcommand{\Exp}{\mathbb{E}}	% expectation
\newcommand{\iid}{\text{i.i.d.}}
\newcommand{\indep}{{\perp\!\!\!\!\perp}}
\newcommand{\noindep}{{\not\!\perp\!\!\!\!\perp}}

\newcommand{\Var}{\operatorname{Var}}	% variance
\newcommand{\Cov}{\operatorname{Cov}}	% covariance
\newcommand{\Corr}{\operatorname{Corr}}

% binary operators
\newcommand{\dist}{\sim}
\newcommand{\distiid}{\overset{\tiny\iid}{\dist}}
\newcommand{\distind}{\overset{\tiny\text{indep}}{\dist}}

% parametric distributions
\newcommand{\distNorm}{\operatorname{N}}	% normal
\newcommand{\distMVN}{\operatorname{MVN}}	% multivariate normal
\title[Dual Induction CLT for high-dimensional m-dependent data: Theorem 3.3]{Dual Induction CLT for High-dimensional $m$-dependent Data: the projected subgradient approximation to the entrywise $\ell_\infty$-projection of the covariance estimator (Theorem 3.3)}
\author{Heejong Bong, Arun Kumar Kuchibhotla and Alessandro Rinaldo}
\date{Source: arXiv:2306.14299v3 (CC BY 4.0)}
\begin{document}
\maketitle
\noindent\emph{Source and licence.} Dual Induction CLT for High-dimensional $m$-dependent Data, arXiv:2306.14299v3, distributed under the Creative Commons Attribution 4.0 International licence, http://creativecommons.org/licenses/by/4.0/, as recorded in the arXiv OAI metadata for this exact version and shown on the abstract page https://arxiv.org/abs/2306.14299v3. Full rights notice: licenses/arxiv-2306.14299-CC-BY-4.0.txt.\par\smallskip
\noindent\textit{Editorial scope.} Counted unit: Theorem 3.3 of the article (source label \textbackslash{}label\{thm:psg\_error\}, source0.tex lines 889--919), the finite-sample bound for the projected subgradient approximation of the entrywise $\ell_\infty$-projection of the covariance estimator, delivered with its complete original proof at source lines 4949--5048 as the single primary proof body. The article is an imsart[aos] article; that publisher class is not part of the pinned runtime, so the unit is delivered inside the AMS amsart wrapper, which carries the article\textquotesingle s own macro and theorem declaration block copied verbatim from its own 3\_preamble.sty, with the single editorial change that the lemma environment is given its own section-based counter so that the article\textquotesingle s Lemma A.13 is also named as a lemma by the cross-reference machinery, together with the editorial counters that reproduce the article\textquotesingle s own numbering. Every printed number of the delivered unit therefore coincides with the numbering the article\textquotesingle s own source assigns: sections 2, 2.1, 2.2 and 3, appendix A, subsection A.6, Theorem 3.3, Lemma A.13, and equations (4), (5), (6), (7), (11), (12), (14), (15), (A.12) and (A.13). Required context retained uncounted: the whole of section 2.1 Setting and Assumptions with the definition of $m$-ring dependence, of $X_I$ and $Y_I$, of the Kolmogorov distance $\mu$, of the conditional-variance constants and of Assumptions (MIN-VAR), (MIN-EV) and (VAR-EV), together with the moment functionals $L_{q,i}$, $\nu_{q,i}$ and their averages; the section 3 paragraphs that define the naive covariance estimator $\hat\Sigma$ and the projection $\tilde\Sigma$; the projected subgradient iteration and the subgradient of the proof\textquotesingle s hypothesis; the appendix heading; the derivation of the anti-concentration bound (A.12); and the statement of Lemma A.13 on which the proof\textquotesingle s probability bound rests. Citations: the article marks its citations with natbib author-year commands (\textbackslash{}citep, \textbackslash{}citet, \textbackslash{}citealp, \textbackslash{}citetlias) and stores its bibliography in the imsart author-year XML dialect, which cannot be transcribed into a standard thebibliography; the citation map of this package is defined on \textbackslash{}cite-form keys only, and the retained spans contain no such key. The two natbib markers inside the counted proof are therefore omitted under the declared bounded\_source\_comparison fidelity receipt, and their published entries, transcribed from the article\textquotesingle s own 0\_aos-all.bbl, are given here so that no attribution is lost: [Boyd, S., Xiao, L. and Mutapcic, A. (2003). Subgradient methods. Lecture notes of EE392o, Stanford University, Autumn Quarter, 2004.] and [Golub, G. H. (1973). Some modified matrix eigenvalue problems. SIAM Review 15.] Accordingly the delivered unit carries no bibliography of its own. All other natbib markers, the article\textquotesingle s numerical studies, its discussion and its remaining appendices lie outside the retained spans. No mathematics is written, repaired or re-derived; the article\textquotesingle s typographical slips are kept literal, including the spacing in "computed in $O(p^{2})$ operations".
\setcounter{section}{1}
\section{High-dimensional Berry--Esseen bound for \texorpdfstring{$m$}{m}-ring dependent random samples with nondegenerate covariance matrices} \label{sec:theoretical_result}

 \subsection{Setting and Assumptions}\label{sec:setting_assumption}

\setcounter{equation}{3}
For any non-empty subset $I \subseteq \{1, 2, \ldots, n\}$, define
\begin{equation*}
    \mathscr{X}_I \equiv \{ X_i: i \in I \}, ~~
    \mathscr{Y}_I \equiv \{ Y_i: i \in I \}.
\end{equation*}
and
\begin{equation} \label{eq:X_I_Y_I}
    X_I \equiv \sum_{i \in I} X_i, ~~
    Y_I \equiv \sum_{i \in I} Y_i.
\end{equation}
% \begin{equation*}
%     \Exp[Y_i Y_j^\top] = \Exp[X_i X_j^\top],\quad\mbox{for all}\quad  i, j \in [1,n].
% \end{equation*}
%{\color{red}\bf \sout{You may want to formally define what $\ints_n$ is.}}
%{\color{red}\bf \sout{You have to make a statement about the joint distribution of $(Y_1, \ldots, Y_n)$ here so that $Y_1, \ldots, Y_n$ are well-defined. With the statements already present, it is not clear how $Y_1, \ldots, Y_n$ are defined.}}
To streamline our discussion, we introduce a special notation for integral intervals. For any pair of integers $i$ and $j$ that satisfy $1 \leq i < j \leq n$, we will denote the ordered index set $\{i, \dots, j\}$ as $[i,j]$. Accordingly, $[1,n]$ will refers to the complete time course. To align with the conventions of real intervals, we will employ parentheses to represent open-ended intervals; e.g., $(1,n]$ denotes the set $\{2, \dots, n\}$. % We also define index set complement: for any index set $I \subset [1,n]$, $I^\cmpl \equiv [1,n] \setminus I$. \commhb{Define those on arbitrary index set $I$.}
With this notation in place, our goal is to derive finite sample bounds for the  quantity
\begin{equation}\label{eq:mu}
\mu\left( X_{[1,n]}, Y_{[1,n]} \right) \equiv \sup_{A \in \mathcal{R}_p} \left| \mathbb{P} \left( \frac{1}{\sqrt{n}} X_{[1,n]} \in A \right) - \mathbb{P} \left( \frac{1}{\sqrt{n}} Y_{[1,n]} \in A \right) \right|.
\end{equation}
To that effect, in order to ensure non-trivial results, we impose appropriate conditions on the conditional variances of the $Y_i$'s which, to the best of our knowledge, are novel. 
Specifically, for any $I, I' \subset [1,n]$, let %${\sigma}_{\min,I}$ and $\underline{\sigma}_{I|I'}$ be
\begin{equation*} \begin{aligned}
    {\sigma}^2_{\min,I} \equiv
    & \min_{k \in [p]} \Var[Y_{I,k}], 
    \quad\mbox{and}\quad
    \underline{\sigma}^2_{I|I'} \equiv
     \lambda_{\min}(\Var[Y_{I}|\mathscr{Y}_{I'}]),
\end{aligned} \end{equation*}
where the superscript $Y_{I,k}$ denotes each $k$-th element of $p$-dimensional random vector $Y_{I}$.
% , and $I^\cmpl$ is the index set complement. i.e., $I^\cmpl \equiv [1,n] \setminus I$.
% \commhb{\sout{Define interval}} {\color{red}\bf \sout{You haven't defined what $[i, j]$ and $(j, i)$ are. I think you are defining this using mod, but it would be a good idea to expand the discussion on the notation.}}
We recall that the conditional variance coincides with the Schur complement of the marginal covariance matrix, i.e.
\begin{equation*}
\begin{aligned}
    \Var[Y_{I}|\mathscr{Y}_{I'}] 
    = \Var[Y_{I}]
    - \Cov[Y_{I}, \mathrm{vec}(\mathscr{Y}_{I'})] \Var[\mathrm{vec}(\mathscr{Y}_{I'})]^{-1} \Cov[\mathrm{vec}(\mathscr{Y}_{I'}), Y_{I}],
\end{aligned}
\end{equation*}
% where $\mathscr{Y}_{(j,i)}$ was considered as the $p\cdot\abs{(j,i)}$-dimensional {\color{red}\bf Why is there $p\cdot|(j,i)|$?} random vector, $(Y_{j+1}^\top, \dots, Y_{i-1}^\top)^\top$. 
where $\mathrm{vec}(\mathscr{Y}_{I'})$ indicates the vectorized representation by concatenation. i.e., $\mathrm{vec}(\mathscr{Y}_{I'}) \equiv (Y_i^\top: i \in I')^\top$. % \commhb{\sout{explicitly say that it is the concatenation of random vectors}}{\color{red}\bf \sout{Why is there $p\cdot|(j,i)|$?}}
In the case $m=0$, i.e. when $X_1,\dots,X_n$ are independent, then 
$\Var[Y_{I}|\mathscr{Y}_{I'}] = \Var[Y_{I}].$
We assume that there exist constants $\sigma_{\min}, \underline{\sigma} > 0$ such that 
for all $I, I' \subset [1,n]$,
\begin{align}
    \sigma^2_{\min,I} & \geq \sigma^2_{\min} \cdot \abs{I}, 
    \label{assmp:min_var} \tag{MIN-VAR} \\
    \underline{\sigma}^2_{I|I'} & \geq \underline{\sigma}^2 \cdot \abs{I \cap I'^{\indep}}, 
    \label{assmp:min_ev} \tag{MIN-EV} \\
    \sigma_{\min} & \le \underline{\sigma} \sqrt{\log(4ep)/2},
    \label{assmp:var_ev}\tag{VAR-EV}
\end{align}
where $\abs{I}$ is the number of elements in $I$, and $I'^{\indep} \;:=\; \bigl\{\,j \in [n] : X_j \;\indep\; \mathscr{X}_{I'}\bigr\}$ is the set of indices whose corresponding variables are independent of $\mathscr{X}_{I'}$. Thus $I \cup I'^{\indep}$ is the subset of $I$ which are separated from $I'$ by at least $m$ positions. 
% \textcolor{red}{Remark that when $m=1$, $\underline{\sigma}^2_{I} = 0$, if $|I|=2$!} \commhb{That's intended: consider the case where $X_{2k} = X_{2k+1}$ for every $k$.} \textcolor{red}{Yes, but we would not have a CLT approximation then, no matter how large $n$ is.} \commhb{We can still have a CLT approximation, because $X_1 \indep X_2 = X_3 \indep X_4 = X_5 \indep \dots$ and so on.} \commhb{Let's leave a remark that the assumption allows this example.} 

It is worth highlighting that Assumption~\eqref{assmp:min_ev} accommodates the possibility of complete dependence between $X_{I}$ and $\mathscr{X}_{I'}$. % adjacent elements, for $\abs{I} \leq 2m$. % This occurs when the interval $I$ is sufficiently short, for instance, when $\abs{I} \leq 2m$. 
For example, consider a scenario where $X_{2j-1} = X_{2j}$ for $j \in [\ceil{\frac{n}{2}}]$, while $X_1, X_3, \dots$ are independent. In this case, $\sigma^2_{[2,3]|\{1\}\cup[4,n]} = \Var[Y_{[2,3]} | Y_1, Y_4, Y_5, \dots, Y_n] = 0$ because $Y_2$ and $Y_3$ exhibit total dependence with $Y_1$ and $Y_4$, respectively. However,  $\sum_{i=1}^n X_i$ still  convergences weakly to $\sum_{i=1}^n Y_i$. Our Berry--Esseen bounds apply to this scenario.


We will also consider two different types of moment bounds: a marginal one on the individual coordinates of the $X_i$'s and of the $Y_i$'s, and a stronger one involving their $L_\infty$ norm. Specifically, for $q \geq 1$ and $i \in [n]$, let 
\begin{equation}\label{eq:definition-L-nu}
    L_{q,i} \equiv \max_{k \in [p]} \Exp[\abs{X_{i,k}}^q] + \max_{k \in [p]} \Exp[\abs{Y_{i,k}}^q],\quad\mbox{and}\quad
    % \mathrm{max} \left\{
    % \Exp[\norm{X_{i}}_\infty], \Exp[\norm{Y_{i}}_\infty] \right\},
% \end{equation*}
% % \begin{equation*}
% %     \overline{\sigma}^2 \equiv
% %     \mathrm{max} \left\{
% %     \lambda_{\max}(\Var[X_i]):
% %     i \in [1,n] \right\},
% % \end{equation*}
% \begin{equation*}
    \nu_{q,i} \equiv \Exp[\norm{X_{i}}_\infty^q] + \Exp[\norm{Y_{i}}_\infty^q].
    % \mathrm{max} \left\{
    % \Exp[\norm{X_{i}}_\infty^3], \Exp[\norm{Y_{i}}_\infty^3]\right\}.
\end{equation}
Clearly, $L_{q,i} \leq \nu_{q,i}$ for all $q$ and $i$.
%Denote the averages of $L_{q,i}$'s and $\nu_{q,i}$'s by
%\begin{equation}\label{eq:definition-L-bar-and-nu-bar}
%    \bar{L}_q = \frac{1}{n} \sum_{i=1}^n L_{q,i} \textand
%    \bar{\nu}_q = \frac{1}{n} \sum_{i=1}^n \nu_{q,i}.
%\end{equation}
For any subset $I \subset [1,n]$, we denote the average marginal and joint moments over the indices in $I$ with
\begin{equation}\label{eq:definition-L-bar-I-and-nu-bar-I}
    \bar{L}_{q,I} \equiv \frac{1}{\abs{I}} \sum_{i \in I} L_{q,i}
    \textand
    \bar\nu_{q,I} \equiv \frac{1}{\abs{I}} \sum_{i \in I} \nu_{q,i}.
\end{equation}
and write $\bar{L}_q \equiv \bar{L}_{q,[1,n]}$ and $\bar{\nu}_q \equiv \bar{\nu}_{q,[1,n]}$ for the global average moments.
We note that, due to Jensen's inequality, $\max_{i \in [n]} \nu_{2,i} \leq \nu_{q,i}^{2/q}$  and $\bar\nu_2 \leq \bar\nu_q^{2/q}$ (since  $q > 2$). 
%%%%


\subsection{Main Results}\label{sec:main}

With the notations and assumptions in place, we are now ready to present our first bound on the quantity $\mu\left( X_{[1,n]}, Y_{[1,n]} \right)$ (see~equation \ref{eq:mu}),  assuming only finite third moments. 

\section{\texorpdfstring{$m$}{m}-dependent bootstrap} \label{sec:m_dep_bootstrap}

\setcounter{equation}{10}
To motivate the proposed procedure, we begin with the simplest case of $m=1$. In this case a natural unbiased estimator for the covariance matrix $\Sigma$ of $Y_{[1,n]}$ is 
\begin{equation} \label{eq:sigma_hat}
    \hat{\Sigma} 
    \equiv \sum_{i=1}^n \left( X_i X_i^\top 
    + X_i X_{i+1}^\top + X_{i+1} X_i^\top \right).
\end{equation}
A straightforward approach is to replace the infeasible distribution of $Y_{[1,n]}$ by the feasible multi-variate Gaussian distribution with mean $\vec{0}$ and covariance $\hat\Sigma$.
However, %a complication arises when sampling from the distribution, since 
$\hat\Sigma$ is not guaranteed to be positive semidefinite, and hence, this Gaussian distribution is not well-defined. 

We begin by 
% constructing the corresponding correlation matrix:
% \begin{equation*}
%     \hat\rho \equiv \hat{D}^{-1/2} \hat\Sigma \hat{D}^{-1/2},
% \end{equation*}
% where $\hat{D} \equiv \operatorname{diag}(\hat\Sigma_{11}, \dots, \hat\Sigma_{pp})$ is the diagonal matrix of marginal variances.  
% We then 
computing the closest positive semidefinite matrix to $\hat\Sigma$ in the elementwise $\ell_\infty$ norm:
\begin{equation} \label{eq:rho_tilde}
    \tilde\Sigma 
    \equiv \underset{\Sigma' \in \mathbb{S}_{+}^p}{\arg\min} \;\norm{\Sigma' - \hat\Sigma}_\infty,
\end{equation}

\setcounter{equation}{13}
\medskip
\textbf{Solving the optimization problem in \cref{eq:rho_tilde}. }
% 
We observe that the optimization problem in \cref{eq:rho_tilde} is convex and can be solved numerically with projected subgradient descent. Starting from $\tilde\Sigma^{(0)} \equiv \diag(\hat\Sigma_{11}, \dots, \hat\Sigma_{pp})$, we iterate
\begin{equation} \label{eq:psg_update}
    \tilde\Sigma^{(k+1)} \equiv \Pi_{\mathbb{S}_{+}^p}(\tilde\Sigma^{(k)} - \eta g^{(k)}),
\end{equation}
where $\eta > 0$ is a fixed step size, $g^{(k)} \in \partial\norm{\Sigma' - \hat\Sigma}_\infty |_{\Sigma' = \tilde\Sigma^{(k)}}$, and $\Pi_{\mathbb{S}_{+}^p}$ denotes projection onto the cone of $p \times p$ positive semidefinite matrices under the Frobenius norm (implemented by eigenvalue thresholding).  A convenient subgradient is obtained by choosing one index that attains the current maximal entry-wise deviation.  Let  
\begin{equation*}
    \mathcal{M}(k)
    \equiv {\arg\max}_{(i,j) \in [p] \times [p]} \norm*{\tilde\Sigma^{(k)}_{ij}-\hat\Sigma_{ij}}
\end{equation*}
and, for any $(i^*,j^*) \in \mathcal M(k)$, let
\begin{equation} \label{eq:subgradient}
    g^{(k)}_{ij}
    = \begin{cases}
      \mathrm{sgn}(\tilde\Sigma^{(k)}_{i^*j^*} - \hat\Sigma_{i^*j^*}),
      & (i,j)=(i^*,j^*) ~\text{or}~ (i,j)=(j^*,i^*),\\
      0, & \text{otherwise}.
    \end{cases}
\end{equation}
The approximation by projected subgradient descent satisfies the following finite sample bound. 
% With the fixed step size $\eta = \norm{\hat\Sigma}_\infty/(p \sqrt{n})$ the following bound holds.


\setcounter{theorem}{2}
\begin{theorem} \label{thm:psg_error}
    Let $\{\tilde\Sigma^{(k)} : k \in [K]\}$ be the sequence generated by \cref{eq:psg_update} with $K \in \nats$ and $\eta = p \norm{\hat\Sigma}_\infty/\sqrt{K}$, and define $k^* \equiv {\arg\min}_{0 \leq k \leq K} \norm{\tilde\Sigma^{(k)}-\hat\Sigma}_\infty$.
    % Suppose that 
    % \begin{equation*}
    %     \frac{\bar{\nu}_{q}^{2/q}}{\norm{\Sigma}_\infty}  \left(\frac{\log(2p)}{n}\right)^{1-2/q} \rightarrow 0.
    % \end{equation*} 
    % If $q \geq 4$, suppose additionally that 
    % \begin{equation*}
    %     \frac{\bar{L}_{4}^{1/2}}{\norm{\Sigma}_\infty} \sqrt{\frac{\log(2p)}{n}} \rightarrow 0.
    % \end{equation*}
    % Assume that $q \geq 3$ and that 
    % \begin{equation*}
    %     \frac{\bar{\nu}_{q}^{2/q}}{\norm{\Sigma}_\infty}  \left(\frac{\log(2p)}{n}\right)^{1-2/q} + \mathbf{1}\{q \geq 4\} \cdot \frac{\bar{L}_{4}^{1/2}}{\norm{\Sigma}_\infty} \sqrt{\frac{\log(2p)}{n}} \rightarrow 0.
    % \end{equation*}
    % {\bf ALE: how about this? Assume that $q \geq 4$ and that that 
    % \begin{equation*}
    %     \frac{\bar{\nu}_{q}^{2/q}}{\norm{\Sigma}_\infty}  \left(\frac{\log(2p)}{n}\right)^{1-2/q} \rightarrow 0 \quad \text{and} \quad  \frac{\bar{L}_{4}^{1/2}}{\norm{\Sigma}_\infty} \sqrt{\frac{\log(2p)}{n}} \rightarrow 0.
    % \end{equation*}}
    Then, with probability at least $1 - \delta'$,
    \begin{equation*}
        \norm{\tilde\Sigma^{(k^*)} - \hat\Sigma}_\infty - \norm{\tilde\Sigma - \hat\Sigma}_\infty 
        \leq C \frac{p \norm{\Sigma}_\infty}{\sqrt{K}},
    \end{equation*}
    for some universal constant $C>0$, where
    \begin{equation*}
        \delta' = C p \cdot \exp\left( - \frac{n \norm{\Sigma}_\infty^2}{C \bar{L}_4} \right) 
        + \left( \frac{C}{\norm{\Sigma}_\infty} \right)^{q/2} \left( \frac{\log(2p/\delta_\circ)}{n} \right)^{q/2 - 1} \bar\nu_q,
    \end{equation*}
    and $\delta_\circ = \left( \frac{C}{\norm{\Sigma}_\infty} \right)^{q/2} \left( \frac{1}{n} \right)^{q/2 - 1} \bar\nu_q$.
    The computational cost of each iteration is $O(p^{2})$, so the overall computational complexity is $O(p^{2} K)$.
\end{theorem}

\appendix
\setcounter{equation}{0}
\renewcommand{\theequation}{\thesection.\arabic{equation}}
\section{Proof of Theorems} \label{sec:pf_thm}

\setcounter{equation}{11}
In other words, $\hat\Sigma$ is centered around $\Sigma$, and an upperbound for $\Delta$ is driven by the law of large numbers. Specifically, for $i \in [n]$, let 
\begin{equation*}
\begin{aligned}
    S_i 
    & \equiv  X_i X_i^\top + X_i X_{i+1}^\top + X_{i+1} X_i^\top
    - \Exp\left[
        X_i X_i^\top + X_i X_{i+1}^\top + X_{i+1} X_i^\top
    \right].
\end{aligned}
\end{equation*}
The summands $(S_i: i \in [n])$ are $3$-ring dependent $p \times p$ mean zero random matrices, and $\Delta = \norm*{\frac{1}{n} \sum_{i=1}^{n} S_i}_\infty$. We divide these summands into three groups: 
for $l \in [3]$ and $j \in [\floor{n/3}]$, let $T^{(l)}_j \equiv S_{3(j-1)+l}$, except for $T^{(3)}_{\floor{n/3}}$, which is defined as $T^{(3)}_{\floor{n/3}} \equiv S_{3\floor{n/3}} + \dots + S_n$. We note that each $(T^{(l)}_j: j \in [\floor{n/3}])$ for $l \in [3]$ is an independent sequence of mean zero Gaussian random matrices, whose entries have finite $q/2$-th moments which are bounded by $3 L_{q,j}$ and $3 \nu_{q,j}$. Based on \cref{thm:conc_ineq_q}, for each $l \in [3]$, we obtain
\begin{equation} \label{eq:conc_ineq_Sigma}
\begin{aligned}
    \norm*{\frac{1}{n} \tsum_{j=1}^{\floor{n/3}} T^{(l)}_j}_\infty 
    & \leq C \bar{L}_{\min\{4, q\}}^{1/2} \left(\frac{\bar{\nu}_{q}}{\delta} \right)^{\max\{2/q-1/2,0\}}
    \left(\frac{\log(2p/\delta)}{n}\right)^{\min\{1-2/q, 1/2\}} \\
    & \quad + C \left(\frac{\bar{\nu}_{q}}{\delta} \right)^{2/q} \left(\frac{\log(2p/\delta)}{n}\right)^{1-2/q},
\end{aligned}
\end{equation}

\setcounter{lemma}{12}
\begin{lemma} \label{thm:conc_ineq_q}
    Suppose that $X_1, \dots, X_n$ are independent $p$-dim random vectors with finite $q$-th moments for some $q > 1$. Specifically, we denote
    \begin{equation*}
        L_{q,i} \equiv \max_{k \in [p]} \Exp[\abs{X_{i,k}}^q],
        \quad
        \nu_{q,i} \equiv \Exp[\norm{X_{i}}_\infty^q],
    \end{equation*}
    $\bar{L}_{q} \equiv \frac{1}{n} \sum_{i \in [n]} L_{q,i}$,
    and $\bar\nu_{q} \equiv \frac{1}{n} \sum_{i \in [n]} \nu_{q,i}$. Then with probability at least $1 - \delta$,
    \begin{equation*}
    \begin{aligned}
        \norm*{\frac{1}{n}\sum_{i=1}^n X_i}_\infty
        & \leq C \bar{L}_{\min\{2, q\}}^{1/2} \left(\frac{\bar{\nu}_q}{\delta} \right)^{\max\{1/q-1/2,0\}}
        \left(\frac{\log(2p/\delta)}{n}\right)^{\min\{1-1/q, 1/2\}} \\
        & \quad + C \left(\frac{\bar{\nu}_q}{\delta} \right)^{1/q} \left(\frac{\log(2p/\delta)}{n}\right)^{1-1/q},
    \end{aligned}
    \end{equation*}
    where $C > 0$ is a universal constant.
\end{lemma}

\setcounter{equation}{12}
\setcounter{subsection}{5}
\subsection{Proof of Theorem 3.3} \label{sec:pf_psg_error}

By well known-properties of sub-gradient methods 
\begin{equation} \label{eq:psg_basic}
    \norm{\tilde\Sigma^{(k^*)} - \hat\Sigma}_\infty - \norm{\tilde\Sigma - \hat\Sigma}_\infty 
    \leq \frac{R^2}{2 \eta K} + \frac{G^2 \eta}{2},
\end{equation}
where $R \equiv \norm{\tilde\Sigma^{(0)} - \tilde\Sigma}_F$, and $G \equiv \max_{k \in [K]} \norm{g^{(k)}}_F$. Due to the choice of $g^{(k)}$ in \cref{eq:subgradient}, $G \leq 2$ almost surely. 
%
% Now we show that $R \leq C \tr(\hat\Sigma)$ with high probability for a universal constant $C > 0$.
We also observe that
\begin{equation*}
    \norm{\tilde\Sigma^{(0)} - \tilde\Sigma}_F
    \leq \norm{\tilde\Sigma^{(0)}}_F + \norm{\tilde\Sigma}_F
    \leq \sqrt{p} \norm{\tilde\Sigma^{(0)}}_\infty + p \norm{\tilde\Sigma}_\infty.
\end{equation*}
% \sout{\bf ALE: where do we use the nuclear norm?} 
Conditional on $\norm{\hat\Sigma - \Sigma}_\infty \leq \frac{1}{2} \norm{\Sigma}_\infty$, 
\begin{equation*}
\begin{aligned}
    \frac{1}{2} \norm{\Sigma}_\infty
    \leq \norm{\Sigma}_\infty - \norm{\hat\Sigma - \Sigma}_\infty
    \leq \norm{\hat\Sigma}_\infty 
    \leq \norm{\Sigma}_\infty + \norm{\hat\Sigma - \Sigma}_\infty
    \leq \frac{3}{2} \norm{\Sigma}_\infty 
\end{aligned}
\end{equation*}
almost surely. At the same time, due to the optimality of $\tilde\Sigma$,
\begin{equation*}
\begin{aligned}
    \norm{\tilde\Sigma}_\infty 
    & \leq \norm{\hat\Sigma}_\infty + \norm{\tilde\Sigma - \hat\Sigma}_\infty \\
    & \leq \norm{\hat\Sigma}_\infty + 2 \norm{\hat\Sigma - \Sigma}_\infty \\
    & \leq \norm{\hat\Sigma}_\infty + \norm{\Sigma}_\infty
    \leq 3 \norm{\hat\Sigma}_\infty,
\end{aligned}
\end{equation*}
almost surely. Hence, under the same condition, plugging in the step size $\eta = p\norm{\hat\Sigma}_\infty / \sqrt{K}$ to \cref{eq:psg_basic},
\begin{equation*}
\begin{aligned}
    \norm{\tilde\Sigma^{(k^*)} - \hat\Sigma}_\infty - \norm{\tilde\Sigma - \hat\Sigma}_\infty 
    & \leq C \cdot \frac{p \norm{\hat\Sigma}_\infty}{\sqrt{K}} \\
    & \leq \frac{3C}{2} \frac{p \norm{\Sigma}_\infty}{\sqrt{K}},
\end{aligned}
\end{equation*}
almost surely for some universal constant $C > 0$.

\smallskip
Now we give an upper bound for probability of the conditioning event, $\norm{\hat\Sigma - \Sigma}_\infty \leq \frac{1}{2} \norm{\Sigma}_\infty$.
According to \cref{eq:conc_ineq_Sigma}, with probability at least $1 - \delta$,
\begin{equation*}
\begin{aligned}
    \norm{\hat\Sigma - \Sigma}_\infty 
    & \leq C \bar{L}_{\min\{4, q\}}^{1/2} \left(\frac{\bar{\nu}_{q}}{\delta} \right)^{\max\{2/q-1/2,0\}}
    \left(\frac{\log(2p/\delta)}{n}\right)^{\min\{1-2/q, 1/2\}} \\
    & \quad + C \left(\frac{\bar{\nu}_{q}}{\delta} \right)^{2/q} \left(\frac{\log(2p/\delta)}{n}\right)^{1-2/q}. \\
\end{aligned}
\end{equation*}
Hence solving 
\begin{equation*}
\begin{aligned}
    \frac{1}{2} \norm{\Sigma}_\infty 
    & \geq C \bar{L}_{\min\{4, q\}}^{1/2} \left(\frac{\bar{\nu}_{q}}{\delta} \right)^{\max\{2/q-1/2,0\}}
    \left(\frac{\log(2p/\delta)}{n}\right)^{\min\{1-2/q, 1/2\}} \\
    & \quad + C \left(\frac{\bar{\nu}_{q}}{\delta} \right)^{2/q} \left(\frac{\log(2p/\delta)}{n}\right)^{1-2/q}, \\
\end{aligned}
\end{equation*}
we obtain
\begin{equation*}
\begin{aligned}
    & \Pr\left[
        \norm{\hat\Sigma - \Sigma}_\infty \geq \frac{1}{2} \norm{\Sigma}_\infty
    \right] \\
    & \leq C p \cdot \exp\left( - \frac{n \norm{\Sigma}_\infty^2}{C \bar{L}_4} \right) 
    + \left( \frac{C}{\norm{\Sigma}_\infty} \right)^{q/2} \left( \frac{\log(2p/\delta_\circ)}{n} \right)^{q/2 - 1} \bar\nu_q,
\end{aligned}
\end{equation*}
where $\delta_\circ = \left( \frac{C}{\norm{\Sigma}_\infty} \right)^{q/2} \left( \frac{1}{n} \right)^{q/2 - 1} \bar\nu_q$ for some universal constant $C > 0$.

% According to the additional assumptions in \cref{thm:psg_error}, 
% \begin{equation*}
%     \Pr\left[
%         \frac{\norm{\hat\Sigma - \Sigma}_\infty}{\norm{\hat\Sigma}_\infty} \leq 1
%     \right] \rightarrow 1
%     ~~ \text{and} ~~
%     \Pr\left[
%         \norm{\hat\Sigma}_\infty \leq 2 \norm{\Sigma}_\infty
%     \right] \rightarrow 1.
% \end{equation*}
% Hence solving the equation of the right hand itaking $\eta = \norm{\hat\Sigma}_\infty/\sqrt{n}$ and $K = p^2 n$, 
% \begin{equation*}
%     \Pr\left[ 
%         \norm{\tilde\Sigma^{(k^*)} - \hat\Sigma}_\infty - \norm{\tilde\Sigma - \hat\Sigma}_\infty 
%         \leq C \frac{\norm{\Sigma}_\infty}{\sqrt{n}}
%     \right] \rightarrow 1,
% \end{equation*}
% for some universal constant $C > 0$. % \sout{\bf ALE: sorry but I don't see this. I get $C (\frac{\norm{\Sigma}_\infty}{p \sqrt{n}} + \frac{\norm{\Sigma}_\infty}{\sqrt{n}})$ instead}

\smallskip
Lastly, each update in \cref{eq:psg_update} has rank at most three, so once the eigendecomposition of $\tilde\Sigma^{(k)}$ is known, the next eigendecomposition can be calculated in $O(p^{2})$ operations.

\end{document}
